\documentclass[11pt]{article}
\usepackage[margin=1in]{geometry}
\usepackage{enumitem}
% =============================================================================
% Preambles/header.tex — shared LaTeX/Beamer preamble for this project
%
% Use from a Beamer lecture by prepending:
%     \input{header}
% and compiling with TEXINPUTS=../Preambles:$TEXINPUTS (the /compile-latex
% skill does this automatically).
%
% PALETTE CONTRACT. Color names defined here MUST match the names in
% Quarto/theme-template.scss so Beamer and Quarto renderings use the same
% palette. The scripts/check-palette-sync.sh script enforces this.
%
% To customize: edit the HEX values below AND the matching $variable in
% Quarto/theme-template.scss, then run scripts/check-palette-sync.sh.
% =============================================================================

% -----------------------------------------------------------------------------
% Engine detection + encoding
%
% Our /compile-latex skill uses XeLaTeX, where inputenc is unnecessary and
% can produce encoding warnings. Only load inputenc under pdfTeX.
% -----------------------------------------------------------------------------
\usepackage{iftex}
\ifPDFTeX
  \usepackage[utf8]{inputenc}
\fi

\usepackage{amsmath, amssymb, amsthm}
\usepackage{booktabs}
\usepackage{graphicx}

% -----------------------------------------------------------------------------
% PALETTE — mirrors Quarto/theme-template.scss variable names.
% Load xcolor and define colors BEFORE anything that references them
% (notably hyperref, hypersetup, and the Beamer color assignments).
% -----------------------------------------------------------------------------
\usepackage{xcolor}

\definecolor{primary-blue}{HTML}{012169}      % headings, accents
\definecolor{primary-gold}{HTML}{B9975B}      % emphasis, borders
\definecolor{highlight-yellow}{HTML}{F2A900}  % markers, alerts
\definecolor{light-bg}{HTML}{E8EDF5}          % title / section backgrounds
\definecolor{jet}{HTML}{1A1A1A}               % body text

% Semantic colors (used in diagrams and annotations)
\definecolor{positive}{HTML}{15803D}          % good / observed / identified
\definecolor{negative}{HTML}{B91C1C}          % bad / problematic / bias
\definecolor{neutral}{HTML}{525252}           % reference / context

% Highlight accents (match .hi-* classes in the SCSS)
\definecolor{hi-slate}{HTML}{314F4F}
\definecolor{hi-green}{HTML}{2E8B57}
\definecolor{hi-red}{HTML}{C41E3A}

% Semantic aliases so skill templates and snippets can write
% \textcolor{accent}{...} without hard-coding "primary-blue".
\colorlet{accent}{primary-blue}
\colorlet{accent2}{primary-gold}

% -----------------------------------------------------------------------------
% Hyperref — loaded AFTER xcolor + palette so link colors resolve.
% -----------------------------------------------------------------------------
\usepackage{hyperref}
\hypersetup{colorlinks=true, linkcolor=primary-blue, urlcolor=primary-blue,
            citecolor=primary-blue, pdfborder={0 0 0}}

% -----------------------------------------------------------------------------
% Course code — set once per deck (after \input{header}, before \begin{document})
% via \coursecode{CS401}. Shown in the footer on every slide so a deck's course
% is identifiable without opening the title page. Empty by default.
% -----------------------------------------------------------------------------
\makeatletter
\newcommand{\coursecode}[1]{\gdef\@coursecodetext{#1}}
\gdef\@coursecodetext{}
\makeatother

% -----------------------------------------------------------------------------
% Beamer theme assignments (only apply under Beamer).
%
% We detect Beamer via \@ifclassloaded{beamer}, which is the documented,
% non-internal way. This must be wrapped in \makeatletter/\makeatother
% because \@ifclassloaded contains an @-catcode character.
% -----------------------------------------------------------------------------
\makeatletter
\@ifclassloaded{beamer}{%
  \usefonttheme{professionalfonts}
  \setbeamercolor{structure}{fg=primary-blue}
  \setbeamercolor{frametitle}{fg=primary-blue}
  \setbeamercolor{title}{fg=primary-blue}
  \setbeamercolor{subtitle}{fg=primary-gold}
  \setbeamercolor{author}{fg=primary-blue}
  \setbeamercolor{institute}{fg=neutral}
  \setbeamercolor{date}{fg=primary-gold}
  \setbeamercolor{itemize item}{fg=primary-blue}
  \setbeamercolor{itemize subitem}{fg=primary-gold}
  \setbeamercolor{alerted text}{fg=primary-gold}
  \setbeamercolor{block title}{fg=white, bg=primary-blue}
  \setbeamercolor{block body}{bg=light-bg}
  % Semantic block variants: block=definition/concept (blue), exampleblock=
  % worked example (green/positive), alertblock=key takeaway or caution (gold).
  % Keep to <=2 colored boxes per slide across ALL three variants (INV-7).
  \setbeamercolor{block title example}{fg=white, bg=positive}
  \setbeamercolor{block body example}{bg=light-bg}
  \setbeamercolor{block title alerted}{fg=white, bg=primary-gold}
  \setbeamercolor{block body alerted}{bg=light-bg}

  % Minimal footer: course code (left) + page X / N (right), no navigation clutter
  \setbeamertemplate{navigation symbols}{}
  \setbeamertemplate{footline}{%
    \color{neutral}\footnotesize\hspace{0.7em}\@coursecodetext\hfill
    \insertframenumber/\inserttotalframenumber\hspace{0.7em}\vspace{0.3em}}
}{}
\makeatother

% -----------------------------------------------------------------------------
% TikZ setup — libraries the snippet gallery uses
% -----------------------------------------------------------------------------
\usepackage{tikz}
\usetikzlibrary{arrows.meta, positioning, calc, decorations.pathreplacing,
                fit, shapes.geometric, backgrounds}

% Shared TikZ styles — snippets and hand-written diagrams can pick these up
% via `\begin{tikzpicture}[every node/.style={...}]` or by naming specific
% styles (e.g., `\node[dag-node] (X) at (0,0) {X};`).
\tikzset{
  % Node styles for causal DAGs and similar
  dag-node/.style = {
    draw=primary-blue, fill=light-bg, rounded corners=2pt,
    minimum width=1.2cm, minimum height=0.9cm, align=center, font=\small
  },
  decision-node/.style = {
    draw=primary-gold, fill=white, diamond, aspect=1.8,
    minimum width=1.8cm, minimum height=1.2cm, align=center, font=\small,
    inner sep=1pt
  },
  % Edge styles — observed vs counterfactual
  observed-edge/.style = {-{Latex[length=2.5mm]}, thick, draw=jet},
  counterfactual-edge/.style = {-{Latex[length=2.5mm]}, thick, draw=neutral, dashed},
  confound-edge/.style = {-{Latex[length=2.5mm]}, thick, draw=negative, dashed},
  % Data-point styles for plots
  observed-dot/.style = {fill=primary-blue, draw=primary-blue, circle, inner sep=1.2pt},
  counterfactual-dot/.style = {fill=white, draw=neutral, circle, inner sep=1.2pt}
}

% -----------------------------------------------------------------------------
% Convenience macros
% -----------------------------------------------------------------------------
\newcommand{\muted}[1]{\textcolor{neutral}{#1}}
\newcommand{\key}[1]{\textcolor{primary-gold}{\textbf{#1}}}
\newcommand{\good}[1]{\textcolor{positive}{#1}}
\newcommand{\bad}[1]{\textcolor{negative}{#1}}

% Standout frame macro: full-bleed dark background for section transitions.
% Beamer-only (uses \begin{frame}/\setbeamercolor/\usebeamerfont) -- do not
% call from an article-class file (e.g. Notes/<CODE>/*-notes.tex). \muted,
% \key, \good, \bad, and \coursecode above are plain xcolor/gdef and are
% safe to use from either class; verified when Notes/article-class support
% was added.
% Expand: \transitionslide{Your transition title}
\newcommand{\transitionslide}[1]{%
  \begingroup
  \setbeamercolor{background canvas}{bg=primary-blue}
  \begin{frame}[plain]
    \centering
    \vfill
    {\usebeamerfont{title}\color{white}\Large #1\par}
    \vfill
  \end{frame}
  \endgroup
}

% Section divider frame (Beamer-only), an alternative to \transitionslide with a
% darker two-line style: near-black (jet) background, a small white label line
% above a larger gold title, and a thin gold rule along the bottom edge.
% Expand: \sectiondivider{Section 2}{Pointers}
\newcommand{\sectiondivider}[2]{%
  \begingroup
  \setbeamercolor{background canvas}{bg=jet}
  \begin{frame}[plain]
    \vspace*{3.0cm}
    {\color{white}\ttfamily\large #1\par}
    \vspace{0.5cm}
    {\color{primary-gold}\LARGE #2\par}
    \begin{tikzpicture}[remember picture, overlay]
      \draw[primary-gold, line width=2.5pt]
        ($(current page.south west)+(0,0.1)$) -- ($(current page.south east)+(0,0.1)$);
    \end{tikzpicture}
  \end{frame}
  \endgroup
}

% -----------------------------------------------------------------------------
% Channel watermark — "Concepts That Click" (YouTube).
%
% Article-class files (CompetitiveExam Q/A sets, Notes, Assignments) get a
% light diagonal draft watermark on every page plus a centered footer naming
% the channel. Beamer decks get a subtle TikZ background watermark instead
% (the footline already carries the course code). Centralized here so every
% MCQ PDF is branded without per-file edits.
% -----------------------------------------------------------------------------
\makeatletter
\@ifclassloaded{article}{%
  \usepackage{draftwatermark}
  \SetWatermarkText{Concepts That Click}
  \SetWatermarkScale{1}
  \SetWatermarkFontSize{2cm}
  \SetWatermarkLightness{0.86}
  \SetWatermarkAngle{45}
  \usepackage{fancyhdr}
  \pagestyle{fancy}
  \fancyhf{}
  \fancyfoot[C]{\footnotesize\textcolor{neutral}{Concepts That Click~|~YouTube: \href{https://www.youtube.com/@concepts.thatclick}{@concepts.thatclick}}}
  \renewcommand{\headrulewidth}{0pt}
  \renewcommand{\footrulewidth}{0pt}
  \fancypagestyle{plain}{%
    \fancyhf{}%
    \fancyfoot[C]{\footnotesize\textcolor{neutral}{Concepts That Click~|~YouTube: \href{https://www.youtube.com/@concepts.thatclick}{@concepts.thatclick}}}%
    \renewcommand{\headrulewidth}{0pt}%
    \renewcommand{\footrulewidth}{0pt}%
  }
}{}
\@ifclassloaded{beamer}{%
  \setbeamertemplate{background}{%
    \begin{tikzpicture}[remember picture, overlay]
      \node[rotate=45, opacity=0.055, align=center]
        at (current page.center)
        {\Huge\textcolor{neutral}{Concepts That Click}};
    \end{tikzpicture}%
  }
}{}
\makeatother 
\coursecode{JKSSB}
\renewcommand{\thesection}{51.\arabic{section}}
\newtheorem{example}{Example}[section]
\newenvironment{definitionbox}[1]{%
  \par\noindent\textbf{Definition (#1).}\ \itshape
}{\par\vspace{0.3em}}
\title{Programming Foundations, Compiler and Data Structures --- Detailed Lecture Notes}
\author{JKSSB Computer Awareness}
\date{\today}

\begin{document}
\maketitle

\section{Programming languages and their levels}
A computer can act only on instructions given in its own \textbf{machine
language}, a pattern of binary 0s and 1s. Writing whole programs in that form
is slow, error-prone, and tied to one processor. Programmers therefore write
in languages that are easier for humans, and a translator converts the
resulting program into machine language. Languages are placed on a
\textbf{level} scale: low-level languages (machine and assembly) are close to
the hardware, while high-level languages are closer to human language. The
lower the level, the more closely the code reflects the actual machine
instructions.

\begin{definitionbox}{Level of a programming language}
How close a language is to the underlying machine. Machine and assembly
languages are low-level; C, C++, Java, and Python are high-level.
\end{definitionbox}

\begin{center}
\begin{tabular}{p{0.24\textwidth}p{0.20\textwidth}p{0.46\textwidth}}
\toprule
\textbf{Language} & \textbf{Generation} & \textbf{Main feature} \\
\midrule
Machine language & First & Binary 0s and 1s; only language the CPU runs directly \\
Assembly language & Second & Short mnemonics such as ADD and MOV; machine-dependent \\
High-level language & Third & English-like and mathematical; machine-independent \\
\bottomrule
\end{tabular}
\end{center}

\section{Machine language}
\textbf{Machine language} is the only language a processor executes directly.
It is written entirely in binary, and it is \emph{machine-dependent}: each
processor family has its own instruction set, so a program written for one
machine does not run on another. Machine language needs no translation, so it
is fast to run, but it is very difficult to write, read, and debug by hand. It
is also called the \emph{first-generation language}.

\section{Assembly language}
\textbf{Assembly language} replaces binary instruction codes with short
\emph{mnemonics} such as ADD, SUB, and MOV. It is still
\emph{machine-dependent}, because each processor family has its own assembly
language, and it maps almost one-to-one onto machine instructions. Assembly
code cannot run directly; it must be translated into machine language by an
\textbf{assembler}. Assembly language is called the \emph{second-generation
language} and is a low-level language.

A common exam trap is to call assembly language a high-level language. It is
not: assembly language is a low-level, second-generation language because its
statements correspond closely to machine instructions. High-level languages,
by contrast, are machine-independent and use English-like statements.

\section{High-level languages}
\textbf{High-level languages} use English-like words and mathematical
notation. Examples are C, C++, Java, Python, COBOL, and FORTRAN. They are
\emph{machine-independent}: after translation, the same source program can run
on different machines. They are easier to learn, write, read, and maintain
than low-level languages, but they must be translated by a \textbf{compiler}
or an \textbf{interpreter} before the CPU can run them.

\section{Generations of programming languages}
Languages are also grouped into generations, which run parallel to the level
scale above. The first three generations match the three levels exactly.

\begin{center}
\begin{tabular}{p{0.22\textwidth}p{0.26\textwidth}p{0.38\textwidth}}
\toprule
\textbf{Generation} & \textbf{Example} & \textbf{Main feature} \\
\midrule
First (1GL) & Machine language & Binary 0s and 1s; run directly \\
Second (2GL) & Assembly language & Mnemonics; needs an assembler \\
Third (3GL) & C, C++, Java & High-level; needs a compiler or interpreter \\
Fourth (4GL) & SQL & Non-procedural; states what, not how \\
Fifth (5GL) & Logic / AI languages & Problem-solving and constraints \\
\bottomrule
\end{tabular}
\end{center}

For this exam, the high-yield rows are the first three: machine language is
first generation, assembly language is second generation, and a high-level
language such as C or Java is third generation.

\section{Language translators: assembler, compiler, interpreter}
A \textbf{translator} is a program that converts source code written in one
language into another form, usually machine language. There are three common
translators, each matched to a different input language and method.

\begin{definitionbox}{Assembler}
A translator that converts an \emph{assembly-language} program into machine
code.
\end{definitionbox}

\begin{definitionbox}{Compiler}
A translator that converts a \emph{high-level} program into machine code (object
code) as a whole, translating the entire program before it runs.
\end{definitionbox}

\begin{definitionbox}{Interpreter}
A translator that converts and executes a high-level program
\emph{line by line}, without producing a separate object file.
\end{definitionbox}

\begin{center}
\begin{tabular}{p{0.20\textwidth}p{0.36\textwidth}p{0.36\textwidth}}
\toprule
\textbf{Point} & \textbf{Compiler} & \textbf{Interpreter} \\
\midrule
Input & Whole program at once & One statement at a time \\
Output & Object code file & No separate object file \\
Execution speed & Faster & Slower \\
Errors & Reported after full translation & Reported line by line \\
Debugging & Harder & Easier \\
Typical examples & C, C++ & Python, JavaScript \\
\bottomrule
\end{tabular}
\end{center}

The compiler translates the whole program in one pass and produces object
code, so execution is fast, but a single error can delay the whole build and
debugging is harder. The interpreter takes one statement at a time, translates
it, and executes it immediately; execution is slower, but errors are shown
line by line, which makes debugging easier.

\section{Source code, object code, and machine code}
These three terms describe a program at different stages of translation.

\begin{definitionbox}{Source code}
The program as written by the programmer in a high-level or assembly language;
it is human-readable text.
\end{definitionbox}

\begin{definitionbox}{Object code}
The machine-readable, relocatable code produced by a compiler or assembler
from source code; it is not yet linked into a final runnable program.
\end{definitionbox}

\begin{definitionbox}{Machine code}
The final binary instructions, in the processor's own language, that the CPU
executes directly.
\end{definitionbox}

The path is: the programmer writes \textbf{source code}; the translator
converts it into \textbf{object code}; after linking, the result is
\textbf{machine code} that the processor runs. Source code is human-readable;
object code and machine code are machine-readable.

A common exam trap is to ask which of these is written by the programmer. The
answer is \textbf{source code}. Object code and machine code are produced by
the translator, not written by hand, and they are stored as binary patterns
rather than readable text.

\subsection*{How a program is translated}
The translation of a high-level program follows a fixed sequence:
\begin{enumerate}[label=\arabic*.]
  \item The programmer writes the program as \textbf{source code}.
  \item A \textbf{compiler} translates the whole source file and produces an
    \textbf{object file}.
  \item A \textbf{linker} joins the object file with any library code needed.
  \item The result is a \textbf{machine-code} executable that the CPU runs.
\end{enumerate}
An interpreter shortens this path: it takes each statement of the source
program, translates it, and executes it immediately, so no separate object
file is produced.

\section{Algorithm}
An \textbf{algorithm} is a finite, step-by-step procedure for solving a
problem. It has a clear start and end, and every step is unambiguous so that
anyone following it reaches the same result. An algorithm is
language-independent: it can be written in plain English or in pseudocode,
without committing to a particular programming language. A program is the
algorithm expressed in a programming language.

\begin{definitionbox}{Algorithm}
A finite, ordered set of unambiguous steps that solves a problem, with a
defined start and end.
\end{definitionbox}

\begin{example}
Adding two numbers is a small algorithm: (1) start, (2) read values $a$ and
$b$, (3) compute $sum = a + b$, (4) display $sum$, (5) stop. The same steps
can later be written in any programming language.
\end{example}

\section{Flowchart}
A \textbf{flowchart} is a pictorial, diagrammatic representation of an
algorithm. It uses standard symbols joined by arrows to show the order in
which steps are carried out, which makes the flow of control easy to see at a
glance.

\begin{center}
\begin{tabular}{p{0.34\textwidth}p{0.50\textwidth}}
\toprule
\textbf{Symbol} & \textbf{Meaning} \\
\midrule
Oval (terminator) & Start and end of the process \\
Parallelogram & Input or output \\
Rectangle & Process (a computation or action) \\
Diamond (rhombus) & Decision (yes/no branch) \\
Arrow & Flow line showing direction \\
\bottomrule
\end{tabular}
\end{center}

A flowchart for the add-two-numbers algorithm uses these symbols in order: an
oval for start, a parallelogram to read the two inputs, a rectangle to compute
the sum, a parallelogram to display the result, and an oval for stop. When a
step involves a choice, a diamond is used, and two arrows leave it, one for
each outcome.

\section{Debugging and types of errors}
\textbf{Debugging} is the process of finding and correcting errors, called
\emph{bugs}, in a program. Three kinds of error are commonly distinguished.

\begin{definitionbox}{Syntax error}
A violation of the rules of the language, such as a missing semicolon or a
misspelled keyword. It is reported by the compiler.
\end{definitionbox}

\begin{definitionbox}{Logical error}
An error in the logic of the program: it compiles and runs but produces a
wrong result. It is the hardest kind to find.
\end{definitionbox}

\begin{definitionbox}{Runtime error}
An error that occurs while the program is running, such as dividing by zero
or reading a file that does not exist.
\end{definitionbox}

A compiler catches syntax errors but not logical errors. A program with a
logical error compiles and runs, and only the output reveals that something
is wrong.

\section{Object-oriented programming basics}
\textbf{OOP} stands for \textbf{Object-Oriented Programming}. It organises a
program around objects rather than around a sequence of actions. Two central
terms are the class and the object.

\begin{definitionbox}{Class and object}
A \emph{class} is a blueprint or template that defines the data and methods;
an \emph{object} is a single instance created from that class.
\end{definitionbox}

Four features define the object-oriented approach.

\begin{itemize}
  \item \textbf{Encapsulation}: binding the data and the methods that operate
    on it into one unit, with the data hidden from outside access.
  \item \textbf{Inheritance}: a new class acquires the properties and methods
    of an existing class.
  \item \textbf{Polymorphism}: ``one name, many forms'' --- the same operation
    behaves differently for different objects.
  \item \textbf{Abstraction}: hiding the implementation details and showing
    only the essential features.
\end{itemize}

As an illustration, \emph{Student} can be a class that defines the data
(roll number, name, marks) and the methods (enrol, show result). A particular
student such as one with roll number 12 is then an \emph{object} of that class.
The class describes the shape; each object is one concrete example of it.

\section{Stack and LIFO}
A \textbf{stack} is a linear data structure in which insertion and deletion
take place at only one end, called the \emph{top}. It follows the order
\textbf{LIFO}: Last In, First Out. The last element pushed onto the stack is
the first element popped from it.

\begin{definitionbox}{Stack (LIFO)}
A linear data structure that works on the Last In, First Out rule, with
insertion and deletion performed at the top.
\end{definitionbox}

Insertion into a stack is called \textbf{PUSH} and deletion is called
\textbf{POP}. The other end of the stack is the \emph{bottom}, and no
insertion or deletion happens there.

\begin{example}
A stack of files on a desk behaves in exactly this way. The file placed on top
last is the first one picked up next. That is the LIFO rule in everyday form.
\end{example}

\section{Basic data structures}
A data structure is a way of organising data so that it can be used
efficiently. The common structures tested at this level are the array, stack,
queue, linked list, and tree.

\begin{center}
\begin{tabular}{p{0.20\textwidth}p{0.64\textwidth}}
\toprule
\textbf{Structure} & \textbf{Defining property} \\
\midrule
Array & Fixed-size elements of the same type, stored in contiguous memory and accessed by an index \\
Stack & Last In, First Out (LIFO); push and pop take place at the top \\
Queue & First In, First Out (FIFO); insertion at the rear, deletion at the front \\
Linked list & Nodes that hold data and a pointer to the next node; not contiguous and can grow dynamically \\
Tree & A hierarchical structure with a root node and child nodes \\
\bottomrule
\end{tabular}
\end{center}

To identify a structure in an exam item, look for its defining rule: same type,
fixed size, and index access means an \textbf{array}; Last In, First Out means
a \textbf{stack}; First In, First Out (like people in a queue) means a
\textbf{queue}; data-plus-pointer nodes that grow dynamically mean a
\textbf{linked list}; and a root with branches means a \textbf{tree}.

The two most commonly confused structures are the stack and the queue. They
differ only in the order in which elements leave the structure.

\begin{center}
\begin{tabular}{p{0.20\textwidth}p{0.34\textwidth}p{0.34\textwidth}}
\toprule
\textbf{Point} & \textbf{Stack} & \textbf{Queue} \\
\midrule
Order & LIFO (Last In, First Out) & FIFO (First In, First Out) \\
Insertion & Push, at the top & Enqueue, at the rear \\
Deletion & Pop, at the top & Dequeue, at the front \\
Both ends used? & No, one end only & Yes, front and rear \\
\bottomrule
\end{tabular}
\end{center}

\section{Exam retrieval and common confusions}
The following distinctions prevent most errors on this topic.
\begin{enumerate}
  \item A \textbf{compiler} translates the \textbf{whole program at once}; an
    \textbf{interpreter} works \textbf{line by line}.
  \item An \textbf{assembler} translates \textbf{assembly language}; a
    \textbf{compiler} and an \textbf{interpreter} translate
    \textbf{high-level} languages.
  \item \textbf{Source code} is human-readable; \textbf{object code} and
    \textbf{machine code} are machine-readable.
  \item An \textbf{algorithm} is a step-by-step procedure; a
    \textbf{flowchart} is its pictorial representation.
  \item A \textbf{stack} is \textbf{LIFO}; a \textbf{queue} is
    \textbf{FIFO}. Do not swap them.
  \item A \textbf{logical error} is not caught by the compiler; the program
    runs but the result is wrong.
  \item \textbf{OOP} stands for Object-Oriented Programming; a class is a
    blueprint and an object is an instance of it.
\end{enumerate}

\begin{example}
An item asks which translator produces a separate object file before
execution. Trace the definitions: the compiler translates the whole
high-level program and produces object code, so the answer is the compiler. If
the stem instead says ``translates and executes line by line'', the answer is
the interpreter, and if it says ``translates assembly language'', the answer
is the assembler.
\end{example}

\subsection*{Exercises}
\begin{enumerate}[label=\arabic*.]
  \item Distinguish machine language, assembly language, and a high-level
    language by level and by dependence on the machine.
  \item State the difference between an assembler, a compiler, and an
    interpreter.
  \item Differentiate source code, object code, and machine code.
  \item Define an algorithm and explain how a flowchart represents one.
  \item Distinguish a syntax error, a logical error, and a runtime error.
  \item Explain LIFO and name the two operations of a stack.
  \item Identify the data structure: (a) elements of the same type in
    contiguous memory, (b) Last In First Out, (c) a root with child nodes.
\end{enumerate}

\subsection*{Solutions}
\begingroup\small
\begin{enumerate}[label=\arabic*.]
  \item Machine language is the first-generation, binary, machine-dependent
    language the CPU runs directly. Assembly language is the second
    generation, uses mnemonics, and is also machine-dependent. A high-level
    language is third generation, uses English-like statements, and is
    machine-independent.
  \item An assembler translates assembly language into machine code. A
    compiler translates a high-level program as a whole and produces object
    code. An interpreter translates and executes a high-level program line by
    line.
  \item Source code is the program written by the programmer in a high-level
    or assembly language and is human-readable. Object code is the
    machine-readable code produced by the translator and is not yet linked.
    Machine code is the final binary code the CPU executes directly.
  \item An algorithm is a finite, step-by-step, unambiguous procedure to solve
    a problem. A flowchart represents it pictorially using standard symbols
    (oval, parallelogram, rectangle, diamond, arrow) joined by flow lines.
  \item A syntax error breaks the rules of the language and is reported by
    the compiler. A logical error lets the program run but gives a wrong
    result. A runtime error occurs while the program is running.
  \item LIFO means Last In, First Out: the last element pushed is the first
    element popped. The two operations are push (insertion) and pop
    (deletion), both performed at the top.
  \item (a) Array, (b) Stack, (c) Tree.
\end{enumerate}
\endgroup
\end{document}
